\documentclass[sigplan,screen]{acmart}

\AtBeginDocument{%
  \providecommand\BibTeX{{%
    Bib\TeX}}}

\setcopyright{acmlicensed}
\copyrightyear{2026}
\acmYear{2026}
\acmDOI{XXXXXXX.XXXXXXX}
\acmConference[FARM '26]{The 14th ACM SIGPLAN International Workshop on Functional Art, Music, Modelling and Design}{August XX, 2026}{Indianapolis, IN, USA}
\acmISBN{}

\begin{document}

\title{Some Cool Title}

\author{}
\email{}
\orcid{}
\affiliation{%
  \institution{}
  \city{}
  \state{}
  \country{}
}

\renewcommand{\shortauthors}{XX et al.}

\begin{abstract}

\end{abstract}

\begin{CCSXML}
<ccs2012>
   <concept>
       <concept_id>10010405.10010469.10010475</concept_id>
       <concept_desc>Applied computing~Sound and music computing</concept_desc>
       <concept_significance>500</concept_significance>
       </concept>
   <concept>
       <concept_id>10003752.10003790.10003795</concept_id>
       <concept_desc>Theory of computation~Constraint and logic programming</concept_desc>
       <concept_significance>500</concept_significance>
       </concept>
 </ccs2012>
\end{CCSXML}

\ccsdesc[500]{Applied computing~Sound and music computing}
\ccsdesc[500]{Theory of computation~Constraint and logic programming}

\keywords{counterpoint, logic programming}

\maketitle

\section{Introduction}
\label{sec:introduction}
Modern logic programming languages offer a promising foundation for procedural content generation with diverse applications across the arts.
However, a foundation is just that, a starting point from which something more is built.
If  works of art are implemented directly as logic programs, the potential artists are limited to those comfortable programming in avant-garde languages.
Formal logic may be as elegant as a violin and as flexible as a paintbrush, but it is a radically different user interface from both, and developing widely-usable procedural content generation tools requires building widely-understandable user interfaces atop our chosen software foundations.
Only then can we support \emph{casual creators} who use computational creativity tools to find inherent joy rather than achieve expertise.

This paper presents \emph{CounterChoice}, a case study in developing user-facing procedural content generation tools using modern logic programming languages.
CounterChoice is a graphical tool for generating first-species counterpoint compositions: the user enters the first voice of the composition (the \emph{cantus firmus}) and a logic program automatically synthesizes a second voice to complete the composition (the \emph{counterpoint}).
CounterChoice is implemented in the finite-choice logic programming language \emph{Dusa}.
Developed in response to the limitations of answer-set programming (ASP), Dusa is a promising language for constraint-intensive search-based generative problems like counterpoint composition, but Dusa is so new that existing applications are limited.
Counterpoint synthesis is an appropriate case study because the search is large enough to pose a computational challenge and the problem constraints are sophisticated but widely agreed-upon since counterpoint is a standard topic in formal music theory education.																				
Counterpoint and other harmony exercises in general are so clearly constraint-solving problems that one of the authors daydreamed ${\approx}14$ years ago of solving her Harmony homework in Prolog. Years later, logic programming technology has advanced to make this daydream possible.

A single case study holds multiple lessons when viewed through different lenses.
The CounterChoice case study seeks to address the following research questions (RQs):
\begin{itemize}
\item How do the rules of counterpoint translate to the rules of logic? 
\item How much of the tool implementation can readily be performed in Dusa, and how much must be performed in traditional languages?
\item How might the tooling of Dusa and similar languages be revised to simplify the development of user-facing generation tools?
\end{itemize}
Answering RQ1 helps us understand the adequacy of logic for formalizing harmony, i.e., whether music theory relies on rules that cannot be formalized.
Answering RQ2 helps us characterize the extent of productivity gains associated with specialized languages.
Answering RQ3 provides formative feedback to guide continuing development of the Dusa language and promote future application development.

We outline the structure of the paper.
The foundations of counterpoint are presented in Section \ref{sec:background}.
The Dusa implementation of the CounterChoice backend is presented in Section \ref{sec:dusa-implementation}.
The implementation of the frontend as a React web application is presented in Section \ref{sec:ui-implementation}.
Lessons and implications are discussed in Section \ref{sec:discussion}.
Related work is discussed in Section \ref{sec:related-work}.
The paper concludes and outlines future work in Section \ref{sec:conclusion}.

\section{Background}
\label{sec:background}

\section{Dusa Implementation}
\label{sec:dusa-implementation}

\section{UI Implementation}
\label{sec:ui-implementation}

\section{Discussion}
\label{sec:discussion}

\section{Related Work}
\label{sec:related-work}

\section{Conclusion and Future Work}
\label{sec:conclusion}


\bibliographystyle{ACM-Reference-Format}
\bibliography{counterchoice}

\end{document}
\endinput
